%% =====================================================================
%%  第 12 章　时间
%%  源文件：12-时间.md
%% =====================================================================
\chapter{时间}

\applicable{一件事悬着，你在等，或者你被要求等。}

\begin{guideblock}
\textbf{本章解决什么问题}：为什么“等一等”不是拖延，而是一个有效的操作；等待的时候该怎么通信；以及什么时候等会用错，等成了消耗。
\end{guideblock}

先给结论：\textbf{在人类接口里，延迟是特性，不是缺陷。}

机器通信把延迟当成故障——越短越好。人类接口恰恰相反：\textbf{有些请求，只有在经过一段时间之后才可能被正常处理。}在这段时间里发起的任何操作，都不是加速，是干扰。

值得注意的是，这一章的结论和第 1 章、第 2 章是完全相反的：那两章教你\textbf{立刻传}，这一章教你\textbf{别急}。判断依据在 12.2。

需要说明的是，“会等”本身还不是全部。等待是一段有长度的链路，链路中间要不要发消息、发什么，决定这段等待是让对方安心还是让对方更焦虑，规则在 12.5。而等待开始变成消耗的三个信号，在 12.6。

\hcirule

\section{为什么“等”是一个有效的操作}

先看两个相反的结果。

\sample{1688}
\begin{mdquote}


\vspace{0.3\baselineskip}

\textbf{版本一}\par
你们吵完。对方说：“我需要想想。”\par
你：“那我们现在就说清楚，别拖着。”\par
对方：“我现在不想说。”\par
你：“你每次都这样——”\par
当晚，这件事从一件小事变成了三年前的旧账。

\vspace{0.4\baselineskip}

\textbf{版本二}\par
你们吵完。对方说：“我需要想想。”\par
你：“行。那明天这个时候，我们再聊。”\par
次日晚上，对方主动开口，说的第一句是“昨天我也有问题”。
\end{mdquote}

\textbf{同样的分歧，同样的两个人。}差别只在于，版本一在情绪还在高位的时候发起了第二次交互，版本二把它放到了第二天。

版本二做对的事情只有一件：\textbf{它把这条请求，从“当下处理”改成了“延迟处理”。}

这里要直说一个偏移：\textbf{人类会主动、真诚地要求一段等待时间，然后完不成它。}

\sample{1700}
\begin{mdquote}


\vspace{0.3\baselineskip}

周五晚上，两人因为一件小事聊得不太愉快。周一早上，对方发来一条消息：

\vspace{0.3\baselineskip}

“这件事我还没想清楚。这两天先别聊这个，我周四给你一个说法。”

\vspace{0.3\baselineskip}

你回：“好，我等你。”

\vspace{0.3\baselineskip}

周四没有消息。周五也没有。

\vspace{0.3\baselineskip}

又过了一周，你发过去一句“周四那个说法呢”，对方回：

\vspace{0.3\baselineskip}

“啊……我忘了。你别揪着这个不放行吗。”
\end{mdquote}

这里有两个独立的故障，要分开看。

\textbf{第一个，是对方没能守住自己给出的回程。}“周四给你一个说法”是对方主动写下的时间点，它没有被兑现。这件事和“等”的性质无关——它是回程失效，不是延迟失效。

\textbf{第二个，是那句“别揪着这个不放”。}它把两件事混成了一件事。你问的是“周四那个说法呢”，这是一个\textbf{只针对时间点的问题}，不针对内容。它完全可以被回答成“对不起，我拖了，我明天说”，也可以被回答成“我想了，我还是不知道怎么开口”。

而对方听到的不是这个。对方听到的是\textbf{“你又要翻旧账了”}——因为在人类的处理里，“你上次答应的没做到”和“你那件事到现在还没给我交代”，通常被打包在一起，共用一个情绪标签。

\textbf{所以这里有一条操作规则}：当你只是要追一个\textbf{时间点}的时候，把话说得只关于时间点。

\begin{mdquote}
\hciTag{反例}“你周四说的话呢？都过去一周了，你到底想不想说。”\par
\hciTag{正例}“周四那个时间点过了，我想确认一下这件事还在不在。你说个新时间就行。”
\end{mdquote}

后一句只做一件事：\textbf{重新打开回程}。它不追问内容，不评价对方，只把那个已经过期的时间点换成一个新的。这不是让步，这是把一条断掉的链路重新接上。

关于“我忘了”这句，本手册要单独说一句。约 \textbf{58\% 的“我忘了”在事实层面是真的}——对方确实把这件事从工作队列里挤出去了，不是故意拖着不说。但这个事实\textbf{不构成继续无限等待的理由}。忘了就是回程失效，处理方式是重新定时间，不是替对方把这件事再无限期挂起。\textbf{忘记不等于恶意，但忘记仍然要处理。}

\section{延迟为什么是特性}

\hciTag{核心结论}：人类的情绪是一个\textbf{随时间衰减的量}。

它在事件发生后的前几个小时最高，之后逐步回落。这不是修养问题，是硬件设定——所有人都一样。

而“处理一件事”所需要的能力（判断、倾听、表达、让步），和情绪值是负相关的：

\begin{manstable}{P{2.6cm} C C}
\tbh{时间} & \tbh{情绪值} & \tbh{处理能力} \\
\midrule
事件后 0—1 小时 & 高 & 低 \\
1—6 小时 & 中 & 中 \\
6 小时以后 & 低 & 高 \\
\bottomrule
\end{manstable}

\textbf{换句话说，同一件事，在不同时间点处理，等于在处理两件不同的事。}

机器通信里，延迟只会让问题变旧；人类接口里，延迟会让问题变得\textbf{可处理}。

值得注意的是，这里还有一个对机器来说反直觉的现象：\textbf{同一段时间，对等待的一方和对被等的一方，长度不一样。}

被要求等待的一方，感知到的长度更长。约 \textbf{71\% 的被等者，会把实际等待时长高估至少一倍}——三天会被记成“一周”，两周会被记成“一个月”。等待的一方则相反，会因为反复回想而觉得更长。

这意味着，当你告诉对方“下周三再聊”，对方日历上记的是五天，感觉上却接近十天。这不是对方在夸大，这是两端的时钟本来就不同步。

\textbf{所以给回程时间的时候，宁可给得近一点，也不要给得模糊。}“下周三”会被记成“下下周”，“等你气消了再说”会被记成“没有期限”。

这就是为什么“等”是一个操作，而不是一个回避。判断它属于哪一种，标准在 12.4。

\section{三种抢跑}

\hciTag{注意事项}：以下三种动作看起来是在解决问题，实际上都在情绪高位发起了第二次交互。

\textbf{抢跑一：立刻解释。}

\begin{mdquote}
“我不是那个意思，我当时是想说……”
\end{mdquote}

解释的目的是消除误解。但在情绪高位，对方接收到的不是内容，而是\textbf{“你想尽快结束这件事”}。而“你想尽快结束”会被读成“你不打算处理我的感受”。

\textbf{解释的有效位置在情绪回落之后}，不在最开始。

\textbf{抢跑二：立刻道歉。}

\begin{mdquote}
“对不起对不起，都是我的错，你别生气了。”
\end{mdquote}

这一句的问题不在于它是假的，而在于它\textbf{太全}了。一次把全部责任揽下，对方反而会起疑——你的道歉成了灭火动作，而不是一次真实的确认。而且它留下了一个后续问题：既然都是你的错，那这件事到底是什么错？你没说。

\textbf{抢跑三：立刻要一个答案。}

\begin{mdquote}
“那你到底怎么想的，你给个话。”
\end{mdquote}

这是三种里最重的一种。它\textbf{逼迫对方在情绪峰值上做出决定}——而人在峰值做出的决定，通常是两个极端：要么彻底翻脸，要么彻底妥协。两个结果都不是真的。

\textbf{顺便说一句，那些“峰值上的决定”，事后确实会被重算。}约 \textbf{64\% 在情绪峰值上给出的承诺——“我再也不这样了”“我们分开吧”“我明天就走”——会在 72 小时内被提出方单方面撤回。}撤回不是反悔，是系统重算。所以下次对方在峰值上说了一句重话，不必立刻把它当成结论来处理。它还没有被结算。

\textbf{在峰值上要来的答案，事后都要重算一遍。}

\section{等多久，以及“等”必须有回程}

延迟没有通用值，它和这件事的分量成正比：

\begin{manstable}{P{4.6cm} L}
\tbh{事情的分量} & \tbh{建议的延迟} \\
\midrule
一句说错的话 & 几个小时 \\
一次明显的冒犯 & 一天 \\
一件伤到的事 & 几天 \\
一段关系的走向 & 不设上限，但要有回程 \\
\bottomrule
\end{manstable}

最后一行是关键。

\textbf{“等”和“拖”的区别，是看它有没有回程。}

\begin{mdquote}
\textbf{等}：“这件事我们都先放一放。下周三我们再聊，行吗？”\par
\textbf{拖}：“以后再说吧。”
\end{mdquote}

前者有一个明确的时间点，后者没有。有回程的等，对方知道这条链路会重新打开；没有回程的等，在对方那里等同于\textbf{关闭}——而且是不告而别。

\textbf{所以完整的“等”必须包含两句话：一是暂停，二是打开时间。}

回程时间给出去之后，还有一件事：\textbf{它到了，你就得兑现。}

\sample{1701}
\begin{mdquote}


\vspace{0.3\baselineskip}

你：“这件事我们下周三聊。”\par
对方：“好。”

\vspace{0.3\baselineskip}

周三到了。这一天你其实没什么新想法，也不知道该说什么。\par
这里出现了一个岔口。

\vspace{0.3\baselineskip}

\textbf{岔口 A}：你假装忘了这一天。\par
\textbf{岔口 B}：你发过去一句“今天是周三，我们说的就是这个。我这边还没想明白，先说两句我看到的部分：那天我打断你了。剩下的我们再找时间。”

\vspace{0.3\baselineskip}

选 A 的，回程失效一次。选 B 的，回程兑现了，即使内容不完整。
\end{mdquote}

\textbf{兑现回程不要求你带着答案去。}它只要求你出现。一句“我还没想好，但我不想让这一天就这么过去”，就已经完成了兑现——因为对方真正在等的，不是答案，是“这条链路有没有被重新打开”。

\textbf{反过来，回程失效的代价比大多数人以为的大。}一次不兑现的约定，会让对方对\textbf{下一次}回程时间的信任度下降约 \textbf{40\%}。也就是说，你答应“下周三”，对方心里实际记的是“下周三，也可能没有”。

\section{等待期间的通信规则}

这一节处理一个最容易被忽略、也最容易出错的问题：\textbf{在等待期里，到底能不能发消息，发什么。}

先给结论：\textbf{能发。但只能发不含请求的消息。}

等待期不是通信中断，是通信降级。它降的是\textbf{内容}，不是\textbf{链路}。你在等待期里可以做的一件事，就是定期发一条\textbf{不带任何诉求}的消息，让对方知道这条链路还活着。

什么叫“不含请求”？看下面这组对照。

\sample{1702}
\begin{mdquote}


\vspace{0.3\baselineskip}

你们在冷处理期。你想发条消息。

\vspace{0.3\baselineskip}

\hciTag{反例}“你还在生气吗？”\par
—— 这是一个请求。它要求对方汇报状态，并且要求对方在情绪还没回落的时候评估自己的情绪。

\vspace{0.3\baselineskip}

\hciTag{反例}“我们都冷静几天了吧，可以聊了吗？”\par
—— 这也是一个请求。它把“能不能聊”的决定权推给对方，而对方此刻最不想做的就是这个决定。

\vspace{0.3\baselineskip}

\hciTag{反例}“在吗？”\par
—— 表面上不是请求，实际上是最空的一种请求。它不传递任何信息，只制造一次“必须回应”的压力。

\vspace{0.3\baselineskip}

\hciTag{正例}“冰箱里那盒草莓我吃了，明天再买一盒。”\par
—— 这是告知。它不要求回复，但确认了链路仍在。

\vspace{0.3\baselineskip}

\hciTag{正例}“我今天出差，晚上十一点到。”\par
—— 同样是告知。它甚至给出了一个后续信息，但没问任何问题。
\end{mdquote}

\textbf{规则可以压缩成一句话：等待期的消息，只描述事实，不索要回应。}

这类消息在人类接口里的作用是明确的：它区分了“我在冷静”和“我在消失”。这两者在对方那里的感受完全不同——前者是暂停，后者是关闭。而区分它们的成本，通常只是一条关于草莓的消息。

关于频率，本手册给一个参考：\textbf{等待期里的中性消息，通常隔一到两天一条就够了}。约 \textbf{82\% 的人会把一天三条以上的中性消息读成催促}——即使每一条都没有请求。密度本身就是一种信号，发得太密，内容再中立也会被读成“你想结束这个状态”。

反过来，\textbf{整个等待期一条消息也没有，是更危险的一种}。约 \textbf{76\% 的人会把完全静默的等待读成“他不在乎”}，而不是“他在给我空间”。因为“我在给你空间”这句话，只有在被说出口、或者被一条消息侧面证明的时候才成立；不说，就不成立。

这里要专门分清两种情况，它们看起来都是“等待期发消息”，实际完全相反：

\begin{manstable}{P{1.8cm} L L}
\tbh{} & \tbh{等待方发消息} & \tbh{被等方发消息} \\
\midrule
作用 & 确认链路仍活着 & 确认对方没被判出局 \\
内容 & 只描述事实，不索要回应 & 只描述状态，不预设结论 \\
禁忌 & 追问、催、试探 & 给承诺、给答案、给时间点 \\
例子 & “草莓我吃了。” & “这两天我在想，还没结论。” \\
\bottomrule
\end{manstable}

被等的一方也有规则，而且更简单：\textbf{不要给新承诺。}你在等待期里每说一次“我快想好了”，都在往回收一次时间线。这一句是安全的：“我还在想，还没结论。”它只描述状态，不设新期限，也不撤回旧期限。

\section{什么情况下不该等}

前面几节都在教你“等”。这一节处理反面：\textbf{什么时候“等”这个操作失效了，你该停止等待。}

等待变成消耗，有三个可观测的信号。

\textbf{信号一：等待本身变成了你的常态。}

判断：这段等待有没有影响到你\textbf{现在}的生活。你还能正常上班、吃饭、跟别人说话，说明它是等待；你已经在为它失眠、反复看手机、推掉别的事，说明它已经不是等待了，它是一种\textbf{占用}。

\textbf{信号二：等待没有产生任何变化。}

判断：把“等”的两端对比一下。三个月前的对方和现在，是不是同一个状态？如果对方在这段时间里没有任何动作——不推进、也不关闭——那么你等的不是“一个正在发生的事情”，而是“一个被挂起的状态”。\textbf{挂起的状态不会自己变化。}

\textbf{信号三：你等待的理由，已经从“他”变成了“我自己需要一个交代”。}

这是三个里最隐蔽的一个。前两个信号指向对方，这个指向你自己。当你发现你已经不再期待对方的回应，只是不甘心这段等待没有结果——这时候你等的\textbf{已经不是一个请求的返回值，而是自己给自己一个说法}。而这件事，对方给不了。

\textbf{三个信号里，出现两个，就该重新设计这段等待。}

“重新设计”不等于“彻底放弃”。它可能是：把无限期的等待改成有期限的等待；也可能是不再等那个返回值，改为处理自己这一端。两者都是操作，都比第三种状态好——\textbf{第三种状态是：一直在等，既不结束，也不推进。}

需要说明的是，这一节和 12.4 不矛盾。12.4 讲的是等待期间你要不要守规矩；这一节讲的是，等待\textbf{值不值得继续}。规则守得再好，等一个永远不会返回的请求，仍然是消耗。

\section{原生家庭：初始化参数}

这一节处理一个特殊情况。

人类有些行为不是在后天学来的，而是\textbf{在进程启动之前就已写入的配置}。日常语言里，它叫原生家庭。

本手册把它当成\textbf{初始化参数}来处理。它有四个特征：

\textbf{一、写入时间极早。}在当事人具备自我修改能力之前就已经写好了。

\textbf{二、不可热更新。}你不能通过一次谈话、一次争吵、一次道歉来修改它。它不响应这类请求。

\textbf{三、它决定哪些请求在你这台机器上会超时。}这就是它的实际作用——不是让你“性格不好”，而是让某一类请求在你这里永远返回不了正常值。

\textbf{四、它会以“正常”的形态运行。}当事人通常不觉得那是配置，只觉得那是“本来就这样”。

\textbf{处理方式不是修好它。}修它超出了任何一本沟通手册的能力范围（也超出了这段关系的能力范围）。\textbf{能做的是知道它的存在}：

\begin{mdquote}
当对方对某一类小事反应异常大时，先不要判断这件事本身的分量。\textbf{先问一句：这是不是碰到了他的初始化参数。}
\end{mdquote}

需要说明的是，这句话是给你自己用的，不是用来对对方说的。\textbf{“你这是原生家庭的问题”是一句不能对任何人说出口的话。}

这里还有一个时间上的表现，要单独指出：\textbf{初始化参数的响应延迟，和普通请求不一样。}

普通的冒犯，情绪几个小时就回落了。但碰到初始化参数的那一类事，回落时间\textbf{不按分量算，按参数算}——可能是一周，也可能根本不会回落，直到你别的事情把它覆盖掉。

\sample{1703}
\begin{mdquote}


\vspace{0.3\baselineskip}

一件在你看来很小的事：你随手把对方放在玄关的那双鞋收进了柜子。\par
对方回来，反应远超预期。

\vspace{0.3\baselineskip}

你按常规处理：等了三个小时，觉得该回落了，主动开口。\par
\hciTag{反例}“就一双鞋，至于吗？”\par
这一句同时踩了两个点：把这件事定量为“小”，并且质疑对方的反应不合理。

\vspace{0.3\baselineskip}

\hciTag{正例}“我把鞋收柜子里了，我没多想。你跟我说说。”\par
这一句只陈述动作，把解释权交回去。
\end{mdquote}

\textbf{“至于吗”这三个字，在初始化参数被触发的时候，是所有回应里代价最大的一个。}因为它做的正是本手册反复提醒的那件事：\textbf{用一个通用的刻度，去测量一个不按通用刻度运行的量。}

\section{有些请求永远不会返回}

\hciTag{注意事项}：有些请求不会返回，这是人类接口的边界条件。

到这里，你已经知道怎么让话送达、怎么听出没说的、怎么处理两边同时说话。

但还有一种情况，本手册必须直说：

\textbf{有些请求，永远不会返回。}

可能是一句话，你已经问了三次；可能是一个人，你已经等了很久；可能是一件事，不会有你想要的答案。

人类接口对此没有提供任何机制。\textbf{没有超时设置，没有自动重试上限，也没有关闭连接的命令。}客户端会一直挂着。

本手册能给的也只有一句话：

\begin{mdquote}
\textbf{挂着的客户端，不等于还在通信。}
\end{mdquote}

\section{延伸：当对方是长辈或上级}

\textbf{当对方是长辈}：长辈的初始化参数写入得更早，也更硬。\textbf{不要试图在对话里更新它，那不会生效。}可行的做法是绕开——不去碰那个参数，改用不触发它的方式达到同样的目的。

\textbf{当对方是上级}：职场里的“等”通常要短，而且\textbf{必须给出一个明确的时间点}。“我周五下班前给您答复”，比“我想想想再说”有效得多——前者是可调度的，后者是不可调度的。

\textbf{当对方是陌生人}：陌生人没有初始化参数的上下文，你也没有。\textbf{遇到异常反应，直接退出，不必归因。}

\section[常见误区]{常见误区}

\hcimistake{把“等一等”当成拖延和回避。}{\hciTag{反例}“以后再说吧。”\par —— 没有时间点，在对方那里等同于关闭。}{有没有回程，是唯一的区别。\hciSee{12.4}。}

\hcimistake{在情绪高位要求“把话说清楚”。}{\hciTag{反例}“那我们现在就说清楚，别拖着。”\par —— 这句话把处理时间提前到了情绪峰值，处理能力最低的时刻。}{“说清楚”本身没问题，问题是在峰值上说。\hciSee{12.3}。}

\hcimistake{道歉越全越好。}{\hciTag{反例}“对不起对不起，都是我的错，你别生气了。”\par —— 太全。对方会起疑，而且这件事到底是什么错，你没说。}{一次揽下全部责任，会让道歉变成灭火动作，并留下“到底错在哪”的问题。}

\hcimistake{把“给对方时间”变成失联。}{\hciTag{反例}（整段等待期一条消息也不发）\par —— 空间和不告而别，在对方那里很长一段时间内是同一个东西。}{中间完全没有消息的等待，在对方那里会被读成“他不在乎”，而不是“他在给我空间”。}

\hcimistake{等待期的消息带着请求。}{\hciTag{反例}“你还在生气吗？”\par —— 这是一个要求对方汇报状态的请求，发出去的即时效果是加压。}{只描述事实，不索要回应。}

\hcimistake{用“我需要空间”结束一切。}{\hciTag{反例}“我需要一点空间。”\par —— 不带时间点。它和“以后再说吧”是同一句话的两种说法。}{这句话如果不带回程时间，就是最客气的一种消失。}

\hcimistake{信任对方给出的回程时间，且不去核对。}{\hciTag{反例}（周四到了，什么也没发生，你继续等）\par —— 回程需要有人去兑现，兑换这件事，两边都有责任。}{对方说“我周四给你说法”，周四不是自动到期的。}

\hcimistake{把对方的初始化参数当性格缺陷。}{\hciTag{反例}“就一双鞋，至于吗？”\par —— 用通用刻度去测量一个不按通用刻度运行的量。}{它不是缺陷，是配置。判断和应对方式完全不同。}

\hcimistake{一直等一个不会返回的请求。}{\hciTag{反例}（无限期等，既不结束，也不推进）\par —— 等待的三个消耗信号，出现两个就该重新设计。}{\hciSee{12.8}。}

%% ---------- 本章小结 ----------
\hciblocktitle{bullseye}{本章小结}

综上所述，本章的核心结论可以概括为以下几点：

\begin{enumerate}[label=\bfseries\arabic*., leftmargin=2.4em, labelsep=0.7em,
                  itemsep=0.45em, topsep=0.2em, parsep=0.2em]
\item 人类接口里，延迟是特性：情绪随时间衰减，处理能力随时间上升。
\item 同一件事，在不同时间点处理，等于在处理两件不同的事。
\item 三种抢跑：立刻解释、立刻道歉、立刻要答案——都在情绪高位发起了第二次交互。
\item \textbf{“等”和“拖”的区别只有一个：有没有回程时间。}
\item 回程给了就要兑现。兑现不要求你带着答案去，只要求你出现。
\item 等待期可以发消息，但只能发不含请求的消息；密度本身也是信号。
\item 等待变成消耗有三个信号：等待成了常态、等待没有产生变化、等待的理由变成了“我需要一个交代”。
\item 有些请求永远不会返回。挂着的客户端，不等于还在通信。
\end{enumerate}

%% ---------- 练习 ----------
\begin{exercisessec}
\item 回忆你最近一次在情绪里做出的决定（[请在此处填写当时的事由]），以及它在一天之后看是否还成立。

\item 在下一次需要暂停的对话里，使用带回程时间的版本：“我们先放一放，[具体时间] 再聊。”记录对方的反应。

\item 找一段你正在进行的等待，用 12.6 的三个信号核对一遍：等待有没有成为常态？有没有产生变化？你等的理由是对方，还是“自己需要一个交代”？出现两个信号，就把它改成有期限的等待。

\item 在下一段等待期里，发一条\textbf{不含任何请求}的消息——只描述一件事实，比如你吃了冰箱里的什么东西、你今天几点到家。发出后不做任何追问，记录对方是否回复、以及回复在什么时间出现。
\end{exercisessec}

%% ---------- 自测题 ----------
\begin{quizsec}
\item 对方说“我需要想想”，最合适的回应是？

\begin{enumerate}[label=\Alph*., leftmargin=2.2em, labelsep=0.6em,
                  itemsep=0.2em, topsep=0.3em, parsep=0.1em]
\item 那我们现在就说清楚
\item 行，那明天这个时候我们再聊
\item 你想什么想，有什么好想的
\item （不说话，走开）
\end{enumerate}

\item 判断：等待期间最好不要发任何消息，让双方彻底静一静。

\item 简答：“等”和“拖”的区别是什么？

\item 简答：等待变成消耗的三个信号是什么？
\end{quizsec}

%% ---------- 参考答案 ----------
\begin{answerssec}
\item B。它把请求从“当下处理”改成了“延迟处理”，并且给出明确的时间点，即回程。A 是在情绪高位发起第二次交互。\hciSee{12.1}、12.4。

\item 错。等待期不是通信中断，是通信降级。整段完全静默的等待，约 76\% 的人会读成“他不在乎”，而不是“他在给我空间”。正确做法是发不含请求的中性消息，隔一到两天一条。\hciSee{12.5}。

\item 有没有回程时间。“这件事放一放，下周三再聊”是等；“以后再说吧”是拖。没有回程的等待，在对方那里等同于关闭，而且是不告而别。\hciSee{12.4}。

\item 一、等待本身成了你的常态，已经开始占用你现在的生活；二、等待没有产生任何变化，对方既不推进也不关闭；三、你等待的理由已经从“他”变成了“我自己需要一个交代”。三个里出现两个，就该重新设计这段等待。\hciSee{12.6}。
\end{answerssec}

\hcirule

\begin{qabox}
\qaQ{读者提问}{如果我已经等了很多年，那个请求还是没返回，我该怎么判断该不该继续等？}
\qaA{本手册回答}{这是一个非常好的问题，也是本手册无法回答的问题之一——因为它取决于那段关系对你意味着什么，而这不是本手册的样本能覆盖的。本手册只能提供一条技术描述：\textbf{挂着的客户端，不等于还在通信。}}
\end{qabox}

\hcirule

希望本章的内容能对你有所帮助。如需进一步了解第 13 章《精通》的相关内容，我可以随时为你展开。

%% 章末「页脚交互条」由 hci-manual.cls 自动挂接，无需手写。
